\section{The passage of \emph{Les D\'erivateurs}}\label{sec:manuscript}

In Chapter~II of \cite{Der} (manuscript pages 34--36ter; the last of these is on page 23 of
the edition's file for that chapter), Grothendieck shows that filtered
categories are \emph{totalement $W$-asph\'eriques} (Corollaire~2.17; for a
functor $f\colon X\to Y$ the condition is on the category of elements of the
presheaf $y\mapsto\pi_0(y\backslash X)$, which is the condition of
Definition~\ref{def:classes}(c)). He then asks for the converse, for the
class $W_\infty$ of weak equivalences: ``cela implique-t-il que $X$ soit
filtrante~? \c{C}a ne m'\'etonnerait pas'' (ms.~p.~36). A corollary on the
same page states that the opposite $\boldsymbol\Delta^{\circ}$ of the
category of simplices is totally $W_\infty$-aspherical and not filtered. On
pages 36bis--36ter he derives from the argument for it that every connected
simplicial set would be weakly contractible, and concludes: ``Cela me fait
soup\c{c}onner que la conjecture
$X$ totalement $W_\infty$-asph\'erique $\Longrightarrow$ $X$ filtrante
pourrait \^etre valable.''

The results above give the following answers.
\begin{enumerate}
\item The conjecture is false. The category $\Delta_{\mathrm{inj}}$ is exact
 by the theorem of Isbell and Mitchell \cite[Theorem~3]{IM73}, hence total
 by Proposition~\ref{prop:incl}, and it is not filtered. The category $\Cd$
 is a second example, which is not even exact.
\item His doubt about the corollary is right: $\boldsymbol\Delta^{\circ}$ is
 not total. This follows from Proposition~\ref{prop:total}, since the
 action components of $\boldsymbol\Delta^{\circ}$ are the categories of
 simplices of connected simplicial sets; it also follows from
 Proposition~\ref{prop:sifted}, since $\boldsymbol\Delta^{\circ}$ is sifted
 and not filtered.
\item The conjecture holds for sifted categories
 (Proposition~\ref{prop:sifted}), and it holds for categories with a weakly
 terminal object: there even quasi-filtered categories are filtered, by a
 theorem of Isbell and Mitchell \cite[Corollary~4.4]{IM76}.
\end{enumerate}
